% =============================================================================
% CAPÍTULO 10: Reconstrução 3D e Visualização Clínica
% Volume 2 — Pipeline completo: do DICOM segmentado à malha
%            volumétrica pronta para simulação FEM
% Versão: V2 Standalone (reescrito sem \input{})
% Seções:
%   10.1 - Do CBCT à Segmentação
%   10.2 - Pós-Processamento com Open3D
%   10.3 - Geração de Malha Volumétrica para FEM
%   10.4 - Métricas de Validação
%   10.5 - Visualização Clínica e Exportação para Impressão 3D
%   10.6 - Tabela de Ferramentas, Funções e Comandos
% =============================================================================

% --- BADGE DO CAPÍTULO (N2 — Pesquisa Reprodutível) ---
\begin{center}
\begin{minipage}{0.85\textwidth}
\begin{adjustbox}{max width=\textwidth}
\begin{tabular}{|p{0.95\textwidth}|}
\hline
\multicolumn{1}{|c|}{\textbf{\large Badge do Capítulo}} \\
\hline
\textbf{Nível-alvo:} N2 — Pesquisa Reprodutível \\
\textbf{Habilidades:} segmentar CBCT com MONAI, auditar malhas com Open3D,
gerar malha volumétrica com TetGen/Gmsh, calcular métricas Dice e Jaccard \\
\textbf{Pré-requisitos:} Python 3.11+, Open3D~0.18+, MONAI~1.3+, dcm2niix,
Gmsh~4.12+, numpy, scipy \\
\textbf{Comando de verificação:} \verb|tdd_validate.py --chapter 10| \\
\hline
\end{tabular}
\end{adjustbox}
\end{minipage}
\end{center}

\destaque{Nível-alvo N2 — pipeline Open3D para auditoria topológica
de malhas, geração de malha volumétrica e validação por métricas
Dice/Hausdorff, integrado às camadas 2 e 3 do Framework SUS-Twin.}

\label{ch:pipeline}% Se houver conflito ao compilar via light.tex, remova o
% \label{ch:pipeline} correspondente em light.tex (linha 121).

% =============================================================================
\section{Do CBCT à Segmentação}
\label{sec:cbct-segmentacao}
% =============================================================================

A Tomografia Computadorizada de Feixe Cônico (CBCT)\footnote{CBCT
(\textit{Cone-Beam Computed Tomography}) é um tipo de tomografia que
utiliza um feixe de raios X em formato de cone para capturar volumes
tridimensionais da região maxilofacial com baixa dose de radiação.} é
a modalidade de imagem mais difundida na odontologia digital. O
formato DICOM\footnote{DICOM (\textit{Digital Imaging and Communications
in Medicine}) é o padrão internacional para armazenar, visualizar e
transmitir imagens médicas digitais, incluindo tomografias e
radiografias.} resultante precisa ser convertido para um formato
processável por redes neurais — tipicamente NIfTI (\texttt{.nii.gz})
ou metarquivos MHA (\texttt{.mha}).

\subsection{Conversão DICOM $\rightarrow$ NIfTI com dcm2niix}

A ferramenta \texttt{dcm2niix}\footnote{dcm2niix é um conversor de
linha de comando, gratuito e de código aberto, que transforma
arquivos DICOM de scanners de qualquer fabricante para o formato
NIfTI compactado, amplamente utilizado em pipelines de deep
learning.} realiza a conversão preservando metadados essenciais
como tamanho de voxel, orientação e parâmetros de aquisição. O
comando básico é:

\begin{lstlisting}[style=bashstyle, caption={Conversão de DICOM para NIfTI com dcm2niix}, label={lst:dcm2niix}]
dcm2niix -o ./nifti/ -z y -f paciente_001 ./dicom/
# -o: diretorio de saida
# -z y: compactacao gzip
# -f: nome do arquivo de saida
\end{lstlisting}

O resultado são dois arquivos: \texttt{paciente\_001.nii.gz} (volume
de imagem) e \texttt{paciente\_001.json} (metadados no formato
JSON\footnote{JSON (\textit{JavaScript Object Notation}) é um formato
leve de arquivo para armazenar informações estruturadas, como
cabeçalhos de exames, de forma legível tanto por humanos quanto por
máquinas.}).

\subsection{Segmentação Automática com MONAI e DentalSegmentator}

O \textbf{MONAI}\footnote{MONAI (\textit{Medical Open Network for AI})
é um framework de código aberto baseado em PyTorch, desenvolvido
especificamente para deep learning em imagens médicas.} provê a
infraestrutura de carregamento, aumento de dados (\textit{data
augmentation}) e inferência. O \textbf{DentalSegmentator} é um modelo
pré-treinado de segmentação dentomaxilofacial que utiliza uma
arquitetura 3D U-Net com \textit{self-configuring} nnU-Net.

O código abaixo carrega o volume NIfTI, aplica o modelo
pré-treinado e salva a máscara de segmentação:

\begin{lstlisting}[language=Python, caption={Segmentação de CBCT com MONAI e DentalSegmentator}, label={lst:monai_seg}]
import torch
import monai
from monai.transforms import (
    LoadImage, EnsureChannelFirst, ScaleIntensity,
    Resize, Compose
)
from monai.networks.nets import UNet
from monai.inferers import sliding_window_inference
import nibabel as nib
import numpy as np

# Configuracao do dispositivo (GPU ou CPU)
device = torch.device("cuda" if torch.cuda.is_available() else "cpu")
print(f"Dispositivo: {device}")

# Pipeline de transformacao
transform = Compose([
    LoadImage(image_only=True),
    EnsureChannelFirst(),
    ScaleIntensity(minv=0.0, maxv=1.0),
    Resize(spatial_size=(192, 192, 128)),
])

# Carrega e pre-processa o volume
volume = transform("paciente_001.nii.gz")
volume = volume.unsqueeze(0).to(device)  # (1, 1, D, H, W)

# Carrega modelo pre-treinado (DentalSegmentator)
model = UNet(
    spatial_dims=3,
    in_channels=1,
    out_channels=9,  # 8 classes dentarias + fundo
    channels=(16, 32, 64, 128, 256),
    strides=(2, 2, 2, 2),
    num_res_units=2,
).to(device)

# Carrega pesos (simulado: usar torch.load() com checkpoint real)
# model.load_state_dict(torch.load("dental_segmentator.pth"))
model.eval()

# Inferencia por janela deslizante
with torch.no_grad():
    pred = sliding_window_inference(
        volume, roi_size=(96, 96, 96), sw_batch_size=4,
        predictor=model, overlap=0.5
    )
    pred = torch.argmax(pred, dim=1).squeeze(0).cpu().numpy()

# Salva mascara como NIfTI
nifti_mask = nib.Nifti1Image(pred.astype(np.uint8), affine=np.eye(4))
nib.save(nifti_mask, "segmentacao_dental.nii.gz")
print(f"Classes segmentadas: {np.unique(pred)}")
print("Segmentacao concluida.")
\end{lstlisting}

Para uso clínico, o DentalSegmentator\footnote{DentalSegmentator é um
modelo de aprendizado profundo pré-treinado especificamente para
segmentar dentes, mandíbula, maxila e canais mandibulares em
tomografias CBCT.} pré-treinado atinge Dice\footnote{Dice é uma
métrica que mede a sobreposição entre duas segmentações, variando de
0 (nenhuma sobreposição) a 1 (sobreposição perfeita). É a métrica
padrão em competições de segmentação médica como o BraTS.} médio
superior a 0,95 para maxila, 0,93 para mandíbula e 0,89 para dentes
individualizados \cite{dot2024dental}.

% =============================================================================
\section{Pós-Processamento com Open3D}
\label{sec:pos-processamento}
% =============================================================================

A máscara de segmentação volumétrica precisa ser convertida em uma
malha de superfície (\textit{surface mesh}) para posterior geração de
malha volumétrica. O algoritmo de \textbf{Marching Cubes}\footnote{O
algoritmo de Marching Cubes extrai uma superfície poligonal a partir
de um volume voxelizado, criando triângulos que aproximam a fronteira
entre regiões de diferentes intensidades.} extrai a isossuperfície
entre a classe segmentada e o fundo. A malha resultante,
entretanto, contém artefatos que exigem pós-processamento.

\subsection{Suavização Laplaciana e Fechamento de Buracos}

A malha bruta do Marching Cubes apresenta degraus
(\textit{staircase artifacts}) devido à discretização dos voxels. A
suavização Laplaciana\footnote{A suavização Laplaciana desloca cada
vértice da malha para a posição média dos seus vizinhos, reduzindo
rugosidades superficiais sem alterar a topologia global.} iterativa
atenua esses degraus sem colapsar a geometria. Buracos
(\textit{holes}) na malha, comuns em regiões de contato
interproximal, precisam ser fechados para garantir a estanqueidade
(\textit{water-tightness}) necessária à simulação FEM.

\begin{lstlisting}[language=Python, caption={Pós-processamento de malha com Open3D}, label={lst:open3d_post}]
import open3d as o3d
import numpy as np

def postprocess_segmentation(
    mask_path: str = "segmentacao_dental.nii.gz",
    class_id: int = 1,
    voxel_size_mm: float = 0.2,
    smoothing_iters: int = 30,
    output_path: str = "malha_limpa.ply"
) -> o3d.geometry.TriangleMesh:
    """
    Converte mascara segmentada em malha suavizada e estanque.
    
    Args:
        mask_path: Caminho para o arquivo NIfTI da mascara.
        class_id: ID da classe a ser extraida.
        voxel_size_mm: Tamanho do voxel em mm.
        smoothing_iters: Iteracoes de suavizacao Laplaciana.
        output_path: Caminho de saida para a malha PLY.
    
    Returns:
        Malha triangular processada.
    """
    import nibabel as nib
    
    # Carrega mascara segmentada
    nii = nib.load(mask_path)
    mask = nii.get_fdata()
    binary = (mask == class_id).astype(np.uint8)
    print(f"Mascara carregada: {binary.shape}, classe {class_id}")
    
    # Marching Cubes via Open3D (scalar_field)
    mesh, _ = o3d.geometry.TriangleMesh.create_from_point_cloud_alpha_shape(
        o3d.geometry.PointCloud(), alpha=0.0
    )  # Placeholder: usar scikit-image ou o3d voxel grid
    
    # Abordagem com scikit-image (mais robusta)
    from skimage.measure import marching_cubes
    verts, faces, normals, _ = marching_cubes(
        binary, level=0.5, spacing=(voxel_size_mm,) * 3
    )
    
    mesh = o3d.geometry.TriangleMesh()
    mesh.vertices = o3d.utility.Vector3dVector(verts)
    mesh.triangles = o3d.utility.Vector3iVector(faces)
    mesh.compute_vertex_normals()
    
    print(f"Malha inicial: {len(verts)} vertices, {len(faces)} triangulos")
    
    # 1. Suavizacao Laplaciana
    print(f"Aplicando suavizacao Laplaciana ({smoothing_iters} iteracoes)...")
    mesh_smooth = mesh.filter_smooth_laplacian(
        number_of_iterations=smoothing_iters
    )
    mesh_smooth.compute_vertex_normals()
    
    # 2. Fechamento de buracos
    print("Fechando buracos...")
    holes_before = len(mesh_smooth.get_non_manifold_edges())
    mesh_filled = mesh_smooth.fill_holes()
    holes_after = len(mesh_filled.get_non_manifold_edges())
    print(f"Arestas nao-manifold: {holes_before} -> {holes_after}")
    
    # 3. Remocao de componentes desconexas (ilhas)
    print("Removendo componentes isoladas...")
    triangle_clusters, cluster_n_triangles, _ = (
        mesh_filled.cluster_connected_triangles()
    )
    triangle_clusters = np.array(triangle_clusters)
    n_clusters = cluster_n_triangles.shape[0]
    
    # Mantem apenas o maior cluster
    largest_cluster = np.argmax(cluster_n_triangles)
    triangles_to_keep = triangle_clusters == largest_cluster
    mesh_filled.remove_triangles_by_mask(~triangles_to_keep)
    mesh_filled.remove_unreferenced_vertices()
    print(f"Componentes removidas: {n_clusters - 1}")
    
    # 4. Amostragem e simplificacao (se necessario)
    if len(mesh_filled.triangles) > 500_000:
        print(f"Simplificando malha (>{500_000} triangulos)...")
        mesh_simp = mesh_filled.simplify_quadric_decimation(
            target_number_of_triangles=200_000
        )
        mesh_simp.compute_vertex_normals()
        o3d.io.write_triangle_mesh(output_path, mesh_simp)
        print(f"Malha simplificada salva: {output_path}")
        return mesh_simp
    
    o3d.io.write_triangle_mesh(output_path, mesh_filled)
    print(f"Malha processada salva: {output_path}")
    return mesh_filled


if __name__ == "__main__":
    mesh = postprocess_segmentation()
\end{lstlisting}

\subsection{Auditoria Topológica}

Antes de avançar para a geração volumétrica, a malha precisa ser
auditada quanto a:

\begin{itemize}
    \item \textbf{Estanqueidade}: a malha deve ser
        \textit{water-tight}\footnote{Uma malha water-tight é
        geometricamente fechada, sem buracos, formando um volume
        sólido perfeito — condição indispensável para simulações
        de elementos finitos.} — sem buracos na superfície.
    \item \textbf{Manifold}: cada aresta deve pertencer a exatamente
        dois triângulos.
    \item \textbf{Orientação}: todas as normais devem apontar para
        fora do volume.
\end{itemize}

\begin{lstlisting}[language=Python, caption={Auditoria topológica de malha 3D}, label={lst:mesh_audit}]
def audit_mesh(mesh: o3d.geometry.TriangleMesh) -> dict:
    """Audita a qualidade topologica de uma malha triangular."""
    report = {}
    
    # 1. Estanqueidade
    edges = mesh.get_non_manifold_edges()
    report["watertight"] = len(edges) == 0
    report["non_manifold_edges"] = len(edges)
    
    # 2. Vertices e triangulos
    report["n_vertices"] = len(mesh.vertices)
    report["n_triangles"] = len(mesh.triangles)
    
    # 3. Volume e area
    if report["watertight"]:
        report["volume_mm3"] = mesh.get_volume()
        report["area_mm2"] = mesh.get_surface_area()
    else:
        report["volume_mm3"] = None
        report["area_mm2"] = None
    
    # 4. Qualidade dos triangulos (razao aspecto)
    aspects = []
    tris = np.asarray(mesh.triangles)
    verts = np.asarray(mesh.vertices)
    for tri in tris[:1000]:  # amostra de 1000 triangulos
        v0, v1, v2 = verts[tri]
        a = np.linalg.norm(v1 - v0)
        b = np.linalg.norm(v2 - v1)
        c = np.linalg.norm(v2 - v0)
        s = (a + b + c) / 2
        if s > 0:
            area = np.sqrt(s * (s-a) * (s-b) * (s-c))
            aspect = max(a, b, c) / (area ** 0.5 + 1e-8) if area > 0 else 0
            aspects.append(aspect)
    
    report["aspect_ratio_mean"] = np.mean(aspects) if aspects else 0
    report["aspect_ratio_std"] = np.std(aspects) if aspects else 0
    
    # 5. Diagnostico
    report["diagnosis"] = "APROVADA" if report["watertight"] else "REPROVADA"
    if report["watertight"]:
        print(f"Malha APROVADA: {report['n_vertices']} vertices, "
              f"{report['n_triangles']} triangulos, "
              f"volume {report['volume_mm3']:.2f} mm3")
    else:
        print(f"Malha REPROVADA: {report['non_manifold_edges']} "
              f"arestas nao-manifold encontradas")
    
    return report

# Uso:
# mesh = o3d.io.read_triangle_mesh("malha_limpa.ply")
# report = audit_mesh(mesh)
\end{lstlisting}

% =============================================================================
\section{Geração de Malha Volumétrica para FEM}
\label{sec:malha-volumetrica}
% =============================================================================

Com a malha de superfície aprovada na auditoria, o próximo passo é
gerar a malha volumétrica de elementos tetraédricos
(\textit{tetrahedral mesh}) necessária para simulações de Método dos
Elementos Finitos (FEM)\footnote{FEM (\textit{Finite Element Method})
é um método numérico que divide um objeto contínuo em milhares de
elementos pequenos (tetraedros) para calcular tensões, deformações e
outras propriedades físicas através de equações diferenciais.}.

\subsection{Malha Tetraédrica com Gmsh}

O \textbf{Gmsh}\footnote{Gmsh é um gerador de malhas tridimensionais
de código aberto que produz elementos tetraédricos, hexaédricos e
híbridos com controle refinado de densidade local.} é a ferramenta
recomendada por sua flexibilidade e integração com Python via
\texttt{pygmsh} ou API direta.

\begin{lstlisting}[language=Python, caption={Geração de malha volumétrica tetraédrica com Gmsh}, label={lst:gmsh_tet}]
import gmsh
import numpy as np

def generate_tetrahedral_mesh(
    surface_mesh: str = "malha_limpa.ply",
    output_msh: str = "malha_fem.msh",
    lc_max: float = 1.0,
    lc_min: float = 0.1,
    refinement_regions: list = None
):
    """
    Gera malha tetraedrica a partir de malha de superficie.
    
    Args:
        surface_mesh: Caminho para a malha STL/PLY de superficie.
        output_msh: Caminho de saida para o arquivo MSH.
        lc_max: Tamanho maximo dos elementos (mm).
        lc_min: Tamanho minimo dos elementos (mm).
        refinement_regions: Lista de (centro, raio, lc) para refino local.
    """
    gmsh.initialize()
    gmsh.model.add("modelo_dental")
    
    # Importa malha de superficie
    tag = gmsh.model.occ.importShapes(surface_mesh)
    gmsh.model.occ.synchronize()
    
    # Obtem dimensoes para estimar parametros de malha
    volumes = gmsh.model.occ.getEntities(3)
    if not volumes:
        # Cria volume a partir da superficie fechada
        surfaces = gmsh.model.occ.getEntities(2)
        if surfaces:
            surface_loop = gmsh.model.occ.addSurfaceLoop(
                [s[1] for s in surfaces]
            )
            gmsh.model.occ.addVolume([surface_loop])
    gmsh.model.occ.synchronize()
    
    # Configura parametros de malha
    gmsh.option.setNumber("Mesh.MeshSizeMax", lc_max)
    gmsh.option.setNumber("Mesh.MeshSizeMin", lc_min)
    gmsh.option.setNumber("Mesh.Algorithm3D", 1)  # 1=Delaunay, 6=Frontal
    gmsh.option.setNumber("Mesh.Optimize", 1)
    gmsh.option.setNumber("Mesh.OptimizeNetgen", 1)
    
    # Refino local em regioes de interesse (ex: colo do implante)
    if refinement_regions:
        for center, radius, lc in refinement_regions:
            p = gmsh.model.occ.addPoint(center[0], center[1], center[2])
            gmsh.model.occ.synchronize()
            gmsh.model.mesh.field.add("Distance", 1)
            gmsh.model.mesh.field.setNumbers(1, "NodesList", [p])
            gmsh.model.mesh.field.add("Threshold", 2)
            gmsh.model.mesh.field.setNumber(2, "InField", 1)
            gmsh.model.mesh.field.setNumber(2, "SizeMin", lc)
            gmsh.model.mesh.field.setNumber(2, "SizeMax", lc_max)
            gmsh.model.mesh.field.setNumber(2, "DistMin", radius * 0.5)
            gmsh.model.mesh.field.setNumber(2, "DistMax", radius)
            gmsh.model.mesh.field.add("Min", 3)
            gmsh.model.mesh.field.setNumbers(3, "FieldsList", [2])
            gmsh.model.mesh.field.setAsBackgroundMesh(3)
    
    # Gera malha 3D (tetraedros)
    gmsh.model.mesh.generate(3)
    
    # Estatisticas
    element_types, element_tags, _ = gmsh.model.mesh.getElements()
    n_tets = sum(
        len(t) for etype, t in zip(element_types, element_tags)
        if etype == 4  # tetraedro de 4 nos
    )
    n_nodes = len(gmsh.model.mesh.getNodes()[0])
    
    print(f"Malha tetraedrica gerada:")
    print(f"  Tetraedros: {n_tets}")
    print(f"  Nos: {n_nodes}")
    
    # Metricas de qualidade
    qual = gmsh.model.mesh.getQuality()
    print(f"  Qualidade media: {np.mean(qual):.4f}")
    print(f"  Qualidade min:  {np.min(qual):.4f}")
    
    # Salva nos formatos MSH e VTK
    gmsh.write(output_msh)
    gmsh.write(output_msh.replace(".msh", ".vtk"))
    print(f"Malha salva: {output_msh}")
    
    gmsh.finalize()
    return output_msh


if __name__ == "__main__":
    # Exemplo: refino no colo do implante (centro em mm)
    refine = [
        ([0.0, 0.0, 5.0], 2.0, 0.05),  # regiao de contato
    ]
    generate_tetrahedral_mesh(
        surface_mesh="malha_limpa.ply",
        output_msh="malha_fem.msh",
        lc_max=0.5, lc_min=0.05,
        refinement_regions=refine
    )
\end{lstlisting}

\subsection{Parâmetros Críticos da Malha}

A qualidade da malha volumétrica impacta diretamente a precisão e a
convergência da simulação FEM. Três métricas são essenciais:

\begin{table}[htb]
\centering
\caption{Métricas de qualidade da malha volumétrica e seus valores de referência.}
\label{tab:qualidade-malha}
\begin{adjustbox}{max width=\textwidth}
\begin{tabular}{lccp{4cm}}
\toprule
\textbf{Métrica} & \textbf{Aceitável} & \textbf{Ideal} & \textbf{Interpretação} \\
\midrule
\textit{Skewness} & $<0,90$ & $<0,70$ & Distorção angular do elemento \\
\textit{Aspect Ratio} & $<5,0$ & $<3,0$ & Proporção entre maior e menor aresta \\
\textit{Jacobian Ratio} & $>0,3$ & $>0,6$ & Deformação do mapeamento isoparamétrico \\
\textit{Orthogonal Quality} & $>0,1$ & $>0,3$ & Alinhamento com eixos coordenados \\
\textit{Volume Negativo} & 0 & 0 & Elementos invertidos (tolerância zero) \\
\bottomrule
\end{tabular}
\end{adjustbox}
\end{table}

% =============================================================================
\section{Métricas de Validação}
\label{sec:metricas}
% =============================================================================

A validação quantitativa da segmentação e da malha gerada é feita
por três métricas complementares: Dice, Jaccard e distância de
Hausdorff.

\subsection{Coeficiente Dice (F1 Espacial)}

O \textbf{Dice Similarity Coefficient} (DSC)\footnote{O coeficiente
Dice, também conhecido como F1 espacial, mede a sobreposição relativa
entre duas segmentações. Um Dice de 0,95 significa que 95\% dos
voxels coincidem entre a segmentação automática e a manual.} é a
métrica mais utilizada em segmentação médica:

\begin{equation}
\text{DSC}(A,B) = \frac{2|A \cap B|}{|A| + |B|}
\label{eq:dice}
\end{equation}

onde $A$ é a segmentação automática e $B$ é a verdade terrestre
(\textit{ground truth}).

\subsection{Coeficiente Jaccard (IoU)}

O \textbf{Jaccard Index} ou \textbf{Intersection over Union}
(IoU)\footnote{O índice Jaccard mede a razão entre a interseção
(região comum) e a união (região total coberta) de duas
segmentações. É mais conservador que o Dice.} é definido por:

\begin{equation}
J(A,B) = \frac{|A \cap B|}{|A \cup B|}
\label{eq:jaccard}
\end{equation}

Relação entre as métricas: $\text{DSC} = 2J/(1+J)$.

\subsection{Distância de Hausdorff}

A \textbf{distância de Hausdorff}\footnote{A distância de Hausdorff
mede a maior discrepância entre duas superfícies: o máximo da menor
distância de qualquer ponto de uma superfície à outra.} mede a
maior distância entre as superfícies das duas segmentações:

\begin{equation}
H(A,B) = \max\left\{\sup_{a \in \partial A} \inf_{b \in \partial B}
\|a-b\|,\ \sup_{b \in \partial B} \inf_{a \in \partial A}
\|a-b\|\right\}
\label{eq:hausdorff}
\end{equation}

\subsection{Código para Cálculo das Métricas}

\begin{lstlisting}[language=Python, caption={Cálculo de métricas de validação entre segmentações}, label={lst:metrics}]
import numpy as np
from scipy.spatial.distance import directed_hausdorff
import nibabel as nib

def compute_segmentation_metrics(
    pred_path: str = "segmentacao_dental.nii.gz",
    gt_path: str = "ground_truth.nii.gz",
    class_id: int = 1
) -> dict:
    """
    Calcula Dice, Jaccard e Hausdorff entre segmentacao
    predita e ground truth para uma classe especifica.
    """
    # Carrega volumes
    pred_nii = nib.load(pred_path)
    gt_nii = nib.load(gt_path)
    
    pred = (pred_nii.get_fdata() == class_id).astype(np.uint8)
    gt = (gt_nii.get_fdata() == class_id).astype(np.uint8)
    
    # Verifica consistencia dimensional
    assert pred.shape == gt.shape, \
        f"Dimensoes incompativeis: {pred.shape} vs {gt.shape}"
    
    # --- Coeficiente Dice ---
    intersection = np.logical_and(pred, gt).sum()
    pred_sum = pred.sum()
    gt_sum = gt.sum()
    
    dice = (2.0 * intersection) / (pred_sum + gt_sum + 1e-8)
    
    # --- Indice Jaccard (IoU) ---
    union = np.logical_or(pred, gt).sum()
    jaccard = intersection / (union + 1e-8)
    
    # --- Distancia de Hausdorff (95% percentil) ---
    pred_points = np.argwhere(pred)
    gt_points = np.argwhere(gt)
    
    if len(pred_points) == 0 or len(gt_points) == 0:
        hausdorff_95 = float('inf')
    else:
        # Hausdorff bidirecional via scipy
        h1 = directed_hausdorff(pred_points, gt_points)[0]
        h2 = directed_hausdorff(gt_points, pred_points)[0]
        
        # Distancias ponto-a-ponto completas (amostradas)
        from scipy.spatial import cKDTree
        tree_pred = cKDTree(pred_points)
        tree_gt = cKDTree(gt_points)
        
        d_pred_to_gt = tree_pred.query(gt_points)[0]
        d_gt_to_pred = tree_gt.query(pred_points)[0]
        
        # Distancia de Hausdorff 95% percentil
        all_dists = np.concatenate([d_pred_to_gt, d_gt_to_pred])
        hausdorff_95 = np.percentile(all_dists, 95)
        
        # Hausdorff maximo
        hausdorff_max = max(h1, h2)
    
    # --- Sensibilidade e Especificidade ---
    tp = intersection
    fp = pred_sum - intersection
    fn = gt_sum - intersection
    tn = pred.size - (tp + fp + fn)
    
    sensitivity = tp / (tp + fn + 1e-8)
    specificity = tn / (tn + fp + 1e-8)
    precision = tp / (tp + fp + 1e-8)
    f1 = 2 * precision * sensitivity / (precision + sensitivity + 1e-8)
    
    # --- Relatorio ---
    metrics = {
        "class_id": class_id,
        "dice": float(dice),
        "jaccard": float(jaccard),
        "hausdorff_95_mm": float(hausdorff_95),
        "hausdorff_max_mm": float(hausdorff_max) if 'hausdorff_max' in dir() else None,
        "sensitivity": float(sensitivity),
        "specificity": float(specificity),
        "precision": float(precision),
        "f1_score": float(f1),
        "n_voxels_pred": int(pred_sum),
        "n_voxels_gt": int(gt_sum),
    }
    
    return metrics


def print_metrics_report(metrics: dict):
    """Exibe relatorio formatado das metricas."""
    print("=" * 60)
    print(f"RELATORIO DE VALIDACAO - Classe {metrics['class_id']}")
    print("=" * 60)
    print(f"Coeficiente Dice:         {metrics['dice']:.4f}  "
          f"{'EXCELENTE' if metrics['dice'] > 0.90 else 'ACEITAVEL' if metrics['dice'] > 0.80 else 'REVER'}")
    print(f"Indice Jaccard (IoU):     {metrics['jaccard']:.4f}")
    print(f"Hausdorff 95% (voxels):   {metrics['hausdorff_95_mm']:.4f}")
    print(f"Sensibilidade (Recall):   {metrics['sensitivity']:.4f}")
    print(f"Especificidade:           {metrics['specificity']:.4f}")
    print(f"Precisao (Precision):     {metrics['precision']:.4f}")
    print(f"F1-Score:                 {metrics['f1_score']:.4f}")
    print(f"Voxels na predicao:       {metrics['n_voxels_pred']}")
    print(f"Voxels no ground truth:   {metrics['n_voxels_gt']}")
    print("=" * 60)


if __name__ == "__main__":
    metrics = compute_segmentation_metrics(
        pred_path="segmentacao_dental.nii.gz",
        gt_path="ground_truth.nii.gz",
        class_id=1  # mandibula
    )
    print_metrics_report(metrics)
\end{lstlisting}

\subsection{Tabela de Interpretação das Métricas}

\begin{table}[htb]
\centering
\caption{Valores de referência e interpretação clínica das métricas de segmentação.}
\label{tab:metricas-interpretacao}
\begin{adjustbox}{max width=\textwidth}
\begin{tabular}{lccp{5cm}}
\toprule
\textbf{Métrica} & \textbf{Excelente} & \textbf{Aceitável} & \textbf{Interpretação Clínica} \\
\midrule
Dice (estruturas grandes) & $>0,95$ & $0,90$--$0,95$ & Maxila, mandíbula, osso cortical \\
Dice (dentes individuais) & $>0,90$ & $0,80$--$0,90$ & Dentes com esmalte íntegro \\
Dice (LPD) & $>0,70$ & $0,50$--$0,70$ & Ligamento periodontal ($\approx$0,2 mm) \\
Dice (canais) & $>0,80$ & $0,60$--$0,80$ & Canal mandibular, forame mentual \\
Jaccard (IoU) & $>0,85$ & $0,75$--$0,85$ | Equivalente ao Dice -10\% \\
Hausdorff 95\% & $<0,5$ mm & $0,5$--$1,0$ mm & Erro máximo de contorno \\
Sensibilidade & $>0,95$ & $0,85$--$0,95$ & Fração de verdadeiros positivos \\
Especificidade & $>0,95$ & $0,85$--$0,95$ & Fração de verdadeiros negativos \\
\bottomrule
\end{tabular}
\end{adjustbox}
\end{table}

% =============================================================================
\section{Visualização Clínica e Exportação para Impressão 3D}
\label{sec:visualizacao}
% =============================================================================

A etapa final do pipeline é a exportação da malha para formatos
compatíveis com visualizadores clínicos e impressoras 3D. O formato
\textbf{STL}\footnote{STL (\textit{Stereolithography}) é o formato
padrão da indústria para impressão 3D, representando a geometria da
superfície como uma coleção de triângulos.} é o mais universal.

\subsection{Exportação para STL e Visualização}

\begin{lstlisting}[language=Python, caption={Exportação e visualização clínica da malha}, label={lst:stl_export}]
import open3d as o3d
import numpy as np

def export_clinical_visualization(
    mesh_path: str = "malha_fem.msh",
    output_stl: str = "modelo_clinico.stl",
    color_map: dict = None
):
    """
    Exporta malha tetraedrica para visualizacao clinica em STL
    com codificacao de cores por tensao (via mapa de calor).
    
    Args:
        mesh_path: Caminho da malha tetraedrica (.msh ou .vtk).
        output_stl: Caminho de saida STL.
        color_map: Mapeamento {classe: [R, G, B]} para cores.
    """
    import meshio
    
    # Le malha com meshio (suporta MSH, VTK, etc.)
    mesh = meshio.read(mesh_path)
    
    # Se for malha tetraedrica, extrai apenas superficie
    if "tetra" in mesh.cells_dict:
        from meshio import Mesh
        tet_cells = mesh.cells_dict["tetra"]
        
        # Extrai triangulos de superficie dos tetraedros
        # Cada tetraedro tem 4 faces triangulares
        tri_faces = []
        tet_verts = mesh.points
        for tet in tet_cells:
            i, j, k, l = tet
            tri_faces.extend([
                [i, j, k], [i, j, l],
                [i, k, l], [j, k, l]
            ])
        
        # Remove duplicatas (faces internas)
        tri_faces = np.unique(np.sort(tri_faces, axis=1), axis=0)
        
        # Cria malha de superficie Open3D
        tri_mesh = o3d.geometry.TriangleMesh()
        tri_mesh.vertices = o3d.utility.Vector3dVector(tet_verts)
        tri_mesh.triangles = o3d.utility.Vector3iVector(tri_faces)
        tri_mesh.compute_vertex_normals()
    else:
        # Ja e malha de superficie
        tri_mesh = o3d.io.read_triangle_mesh(mesh_path)
    
    # Aplica mapa de cores clinico (ex: por altura)
    verts = np.asarray(tri_mesh.vertices)
    z = verts[:, 2]
    z_norm = (z - z.min()) / (z.max() - z.min() + 1e-8)
    
    # Mapa de calor: azul (base) -> verde (medio) -> vermelho (topo)
    colors = np.zeros((len(verts), 3))
    colors[:, 0] = z_norm       # vermelho
    colors[:, 1] = 1 - z_norm   # verde
    colors[:, 2] = 0.3          # azul fixo
    tri_mesh.vertex_colors = o3d.utility.Vector3dVector(colors)
    
    # Exporta STL colorido (se o formato suportar) ou OBJ
    o3d.io.write_triangle_mesh("modelo_colorido.obj", tri_mesh)
    o3d.io.write_triangle_mesh(output_stl, tri_mesh)
    print(f"Modelo clinico exportado: {output_stl}")
    print(f"Modelo colorido exportado: modelo_colorido.obj")
    
    # Estatisticas para documentacao
    bbox = tri_mesh.get_axis_aligned_bounding_box()
    extent = bbox.get_extent()
    print(f"Dimensoes (mm): {extent[0]:.1f} x {extent[1]:.1f} x {extent[2]:.1f}")
    
    # Visualizacao (opcional, descomentar para uso interativo)
    # o3d.visualization.draw_geometries([tri_mesh])
    
    return output_stl


if __name__ == "__main__":
    export_clinical_visualization(
        mesh_path="malha_fem.msh",
        output_stl="modelo_clinico.stl",
        color_map={1: [0.8, 0.8, 0.8], 2: [0.9, 0.7, 0.5]}
    )
\end{lstlisting}

\subsection{Impressão 3D de Biomodelos}

A malha STL exportada pode ser fisicamente materializada em
impressoras 3D para:

\begin{itemize}
    \item \textbf{Planejamento cirúrgico}: biomodelos\footnote{Biomodelo
        é uma réplica física impressa em 3D da anatomia do paciente,
        usada para planejamento de cirurgias, pré-moldagem de
        instrumentais e comunicação com a equipe cirúrgica.}
        da arcada para pré-moldagem de guias cirúrgicos.
    \item \textbf{Comunicação com o paciente}: modelos táteis para
        explicação do plano de tratamento.
    \item \textbf{Ensino}: réplicas anatômicas para treinamento
        pré-clínico.
\end{itemize}

Parâmetros recomendados para impressão:

\begin{table}[htb]
\centering
\caption{Parâmetros de impressão 3D para biomodelos odontológicos.}
\label{tab:impressao3d}
\begin{adjustbox}{max width=\textwidth}
\begin{tabular}{lcc}
\toprule
\textbf{Parâmetro} & \textbf{Valor} & \textbf{Observação} \\
\midrule
Tecnologia & SLA / DLP & Maior precisão que FDM \\
Resolução de camada & 50 -- 100 $\mu$m & Para detalhes dentários \\
Material & Resina bioinerte & ISO 10993 para contato mucoso \\
Suportes & Automáticos + manuais & Em regiões de socavado (\textit{undercut}) \\
Pós-processamento & Lavagem isopropílica + cura UV & 10 min lavagem + 30 min cura \\
\bottomrule
\end{tabular}
\end{adjustbox}
\end{table}

% =============================================================================
\section{Tabela de Ferramentas, Funções e Comandos}
\label{sec:ferramentas}
% =============================================================================

A \autoref{tab:ferramentas} consolida as principais ferramentas do
pipeline, suas funções e os comandos essenciais para execução.

\begin{table}[htbp]
\centering
\caption{Ferramentas do pipeline de reconstrução 3D: função e comando de uso.}
\label{tab:ferramentas}
\begin{adjustbox}{max width=\textwidth}
\begin{tabular}{llp{4cm}p{5cm}}
\toprule
\textbf{Ferramenta} & \textbf{Função} & \textbf{Comando de Uso} & \textbf{Saída Esperada} \\
\midrule
dcm2niix & Conversão DICOM $\rightarrow$ NIfTI &
\texttt{dcm2niix -o ./nifti/ -z y -f pac\_001 ./dcm/} &
Volume \texttt{.nii.gz} + metadados \texttt{.json} \\
\midrule
MONAI & Framework de deep learning médico &
\texttt{python -c "import monai; print(monai.\_\_version\_\_)"} &
Versão instalada do MONAI \\
\midrule
DentalSegmentator & Segmentação dentomaxilofacial com nnU-Net &
\texttt{python infer.py --input cbct.nii --output seg.nii} &
Máscara NIfTI multiclasse (0--8) \\
\midrule
scikit-image & Marching Cubes (extração de superfície) &
\texttt{skimage.measure.marching\_cubes(volume, level=0.5)} &
Vértices + triângulos + normais \\
\midrule
Open3D & Pós-processamento e auditoria de malhas &
\texttt{open3d.io.read\_triangle\_mesh("malha.ply")} &
Malha carregada para processamento \\
\midrule
Gmsh / pygmsh & Geração de malha tetraédrica volumétrica &
\texttt{gmsh.model.mesh.generate(3)} &
Arquivo \texttt{.msh} com tetraedros \\
\midrule
meshio & Leitura/escrita de malhas (multi-formato) &
\texttt{meshio.read("malha.msh")} &
Objeto Mesh com pontos + células \\
\midrule
TetGen & Geração alternativa de malha tetraédrica &
\texttt{tetgen -pAq1.2 modelo.poly} &
Malha \texttt{.ele} + \texttt{.node} + \texttt{.face} \\
\midrule
CalculiX & Solver FEM de código aberto &
\texttt{ccx\_static -i modelo.inp} &
Arquivo \texttt{.frd} com resultados \\
\midrule
ParaView & Visualização científica pós-processamento &
\texttt{paraview malha\_fem.vtk} &
Interface gráfica interativa 3D \\
\midrule
Blender & Refino manual e texturização de malhas &
\texttt{blender --python export\_mesh.py} &
Malha exportada em formato OBJ/STL \\
\midrule
Ultimaker Cura & Fatiamento para impressão 3D &
\texttt{cura -s modelo\_clinico.stl} &
GCode para impressão FDM/SLA \\
\bottomrule
\end{tabular}
\end{adjustbox}
\end{table}

% =============================================================================
\section{Exercícios Práticos}
\label{sec:exercicios}
% =============================================================================

Os exercícios a seguir consolidam o aprendizado do pipeline completo
de reconstrução 3D. Cada exercício inclui o código de verificação
automática (teste TDD) e o critério de aprovação.

\subsection{Exercício 1: Segmentar CBCT com MONAI}

\textbf{Objetivo}: segmentar um volume CBCT simulado usando MONAI e
extrair a máscara da mandíbula.

\textbf{Instruções}:
\begin{enumerate}
    \item Gere ou carregue um volume CBCT sintético ($64\times64\times64$).
    \item Aplique uma segmentação por limiar de intensidade
        (Hounsfield $>300$ para osso).
    \item Calcule o volume da região segmentada em mm$^3$.
    \item Compare com o \textit{ground truth} fornecido.
\end{enumerate}

\begin{lstlisting}[language=Python, caption={Teste TDD para segmentação CBCT (Exercício 1)}, label={lst:ex1_test}]
import numpy as np
import nibabel as nib

def test_cbct_segmentation():
    """
    Teste de verificacao da segmentacao.
    Criterio: Dice > 0,80 entre segmentacao e ground truth.
    """
    from exercicio1 import segmentar_mandibula
    
    # Gera volume sintetico: esfera de 15 voxels de raio
    shape = (64, 64, 64)
    cx, cy, cz = 32, 32, 32
    r = 15
    x, y, z = np.ogrid[:64, :64, :64]
    gt = ((x - cx)**2 + (y - cy)**2 + (z - cz)**2 <= r**2).astype(np.uint8)
    
    # Simula CBCT com ruido gaussiano
    volume = gt * 500 + np.random.normal(0, 30, shape)
    volume = np.clip(volume, 0, 1000).astype(np.int16)
    
    # Executa segmentacao do aluno
    pred = segmentar_mandibula(volume)
    
    # Calcula Dice
    intersection = np.logical_and(pred, gt).sum()
    dice = 2 * intersection / (pred.sum() + gt.sum() + 1e-8)
    
    print(f"Dice: {dice:.4f}")
    assert dice > 0.80, f"Dice {dice:.4f} < 0.80 — REPROVADO"
    print("TESTE APROVADO: Dice > 0.80")
    return True
\end{lstlisting}

\subsection{Exercício 2: Auditar Malha com Open3D}

\textbf{Objetivo}: carregar uma malha STL, auditar sua qualidade
topológica e reportar métricas de estanqueidade.

\textbf{Instruções}:
\begin{enumerate}
    \item Carregue a malha \texttt{modelo\_teste.stl} com Open3D.
    \item Verifique se a malha é \textit{water-tight}.
    \item Conte vértices, triângulos, arestas não-manifold.
    \item Calcule volume e área da superfície.
    \item Emita diagnóstico ``APROVADA'' ou ``REPROVADA''.
\end{enumerate}

\begin{lstlisting}[language=Python, caption={Teste TDD para auditoria de malha (Exercício 2)}, label={lst:ex2_test}]
import open3d as o3d
import numpy as np

def test_mesh_audit():
    """
    Teste de verificacao da auditoria topologica.
    Criterio: malha deve ser water-tight (0 arestas nao-manifold).
    """
    from exercicio2 import auditar_malha
    
    # Cria malha sintetica: cubo (water-tight)
    mesh = o3d.geometry.TriangleMesh.create_box(width=1, height=1, depth=1)
    o3d.io.write_triangle_mesh("cubo_teste.ply", mesh)
    
    # Executa auditoria
    resultado = auditar_malha("cubo_teste.ply")
    
    print(f"Watertight: {resultado['watertight']}")
    print(f"Vertices: {resultado['n_vertices']}")
    print(f"Triangulos: {resultado['n_triangles']}")
    
    assert resultado['watertight'] == True, "Malha deveria ser water-tight"
    assert resultado['n_vertices'] == 8, "Cubo tem 8 vertices"
    assert resultado['n_triangles'] == 12, "Cubo tem 12 triangulos"
    print("TESTE APROVADO: malha water-tight com topologia correta")
    return True
\end{lstlisting}

\subsection{Exercício 3: Calcular Dice entre Duas Segmentações}

\textbf{Objetivo}: implementar o cálculo do coeficiente Dice entre
duas máscaras NIfTI e interpretar o resultado.

\textbf{Instruções}:
\begin{enumerate}
    \item Carregue duas máscaras NIfTI: \texttt{predicao.nii.gz} e
        \texttt{ground\_truth.nii.gz}.
    \item Calcule Dice, Jaccard e sensibilidade.
    \item Classifique o resultado segundo a \autoref{tab:metricas-interpretacao}.
    \item Salve o relatório em um arquivo \texttt{relatorio\_metricas.json}.
\end{enumerate}

\begin{lstlisting}[language=Python, caption={Teste TDD para cálculo de Dice (Exercício 3)}, label={lst:ex3_test}]
import numpy as np
import json

def test_dice_calculation():
    """
    Teste de verificacao do calculo de Dice.
    Criterio: Dice entre mascaras identicas deve ser 1.0.
    """
    from exercicio3 import calcular_metricas
    
    # Cria duas mascaras identicas (esfera 3D)
    shape = (32, 32, 32)
    cx, cy, cz = 16, 16, 16
    r = 8
    x, y, z = np.ogrid[:32, :32, :32]
    mask = ((x - cx)**2 + (y - cy)**2 + (z - cz)**2 <= r**2).astype(np.uint8)
    
    # Salva como NIfTI temporarios
    import nibabel as nib
    nib.save(nib.Nifti1Image(mask, np.eye(4)), "predicao.nii.gz")
    nib.save(nib.Nifti1Image(mask, np.eye(4)), "ground_truth.nii.gz")
    
    # Calcula metricas
    metrics = calcular_metricas("predicao.nii.gz", "ground_truth.nii.gz")
    
    print(f"Dice: {metrics['dice']:.4f}")
    print(f"Jaccard: {metrics['jaccard']:.4f}")
    
    # Mascaras identicas -> Dice = 1.0, Jaccard = 1.0
    assert abs(metrics['dice'] - 1.0) < 1e-4, \
        f"Dice {metrics['dice']} != 1.0"
    assert abs(metrics['jaccard'] - 1.0) < 1e-4, \
        f"Jaccard {metrics['jaccard']} != 1.0"
    
    # Salva relatorio
    with open("relatorio_metricas.json", "w") as f:
        json.dump(metrics, f, indent=2)
    print("Relatorio salvo: relatorio_metricas.json")
    print("TESTE APROVADO: Dice = Jaccard = 1.0 para mascaras identicas")
    return True
\end{lstlisting}

\subsection{Gabarito dos Exercícios}

O gabarito completo está disponível no repositório do OpenCode
Ecosystem em \texttt{opencode/gabaritos/cap10/}. Para verificar
automaticamente todos os exercícios:

\begin{lstlisting}[style=bashstyle, caption={Verificação automática dos exercícios do Capítulo 10}, label={lst:run_tests}]
# Executa todos os testes do capitulo 10
tdd_validate.py --chapter 10

# Executa teste individual
pytest test_exercicio1.py -v
pytest test_exercicio2.py -v
pytest test_exercicio3.py -v
\end{lstlisting}

% =============================================================================
\section*{Resumo do Capítulo}
\label{sec:resumo}
% =============================================================================

Este capítulo apresentou o pipeline completo de reconstrução 3D e
visualização clínica, desde a conversão DICOM até a malha
volumétrica pronta para simulação FEM. Os principais pontos são:

\begin{itemize}
    \item A conversão DICOM $\rightarrow$ NIfTI com dcm2niix é o
        primeiro passo essencial para processamento computacional.
    \item O MONAI com DentalSegmentator automatiza a segmentação
        das estruturas dentomaxilofaciais com alta precisão (Dice
        $>0,90$ para estruturas grandes).
    \item O pós-processamento com Open3D (suavização, fechamento de
        buracos, remoção de componentes) é indispensável para obter
        malhas \textit{water-tight}.
    \item A geração de malha tetraédrica com Gmsh produz o
        insumo necessário para simulações de elementos finitos.
    \item As métricas Dice, Jaccard e Hausdorff fornecem validação
        quantitativa obrigatória para publicações científicas Qualis A1.
    \item A exportação para STL viabiliza a impressão 3D de
        biomodelos para planejamento cirúrgico e comunicação com
        o paciente.
\end{itemize}

% =============================================================================
% FIM DO CAPÍTULO 10
% =============================================================================
